\documentclass[11pt]{article}

\usepackage{amsmath, amssymb, amsthm, mathtools}
\usepackage[margin=1in]{geometry}
\usepackage{hyperref}
\usepackage{natbib}

\hypersetup{
  colorlinks = true,
  linkcolor  = black,
  citecolor  = black,
  urlcolor   = blue
}

\title{Reinforcing Recursive Language Models: A Literature Review}
\author{Pablo Leyva \\ \small \texttt{first-principles-to-llms} chain, study notes}
\date{\today}

\begin{document}
\maketitle

\iffalse
Stage 5 lit review for the alphaXiv blog "Reinforcing Recursive Language Models"
(NovaSky-AI / SkyRL contributors). Cross-references the chain repo and prior
RLM / GRPO / rubric-reward / options literature. Compiles standalone with
pdflatex + bibtex + pdflatex + pdflatex.
\fi

\section*{Notation}

We adopt the symbol set used in the companion deep dive
(\texttt{02-math-deep-dive.md}). Throughout, $\pi_\theta$ denotes a single shared
autoregressive policy that plays both root and child roles in a recursive
language-model (RLM) trajectory tree.

\begin{center}
\begin{tabular}{ll}
\hline
Symbol & Meaning \\
\hline
$\pi_\theta(a \mid s)$ & Token-level policy: probability of token $a$ given context $s$. \\
$G$                    & Number of root rollouts per query (GRPO group size). \\
$k_g$                  & Number of child rollouts spawned by root $g$. \\
$A_g$                  & Group-relative advantage: $A_g = (r_g - \mu_r)/(\sigma_r + \delta)$. \\
$\rho_\theta(y^{(t)})$ & Importance ratio at token $t$, $\pi_\theta / \pi_{\theta_{\text{old}}}$. \\
$\varepsilon$          & PPO clipping radius (typically $0.2$). \\
$\mathcal{T}$          & Full trajectory tree (root plus descendant rollouts). \\
\hline
\end{tabular}
\end{center}

\section*{Background}

A \emph{recursive language model} (RLM) \citep{zhang2026rlm} treats each LM
session as a process that runs inside a Python REPL sandbox. The sandbox exposes
three privileged calls: \texttt{FINAL(answer)} terminates the rollout,
\texttt{rlm\_query(prompt, context)} spawns a single child rollout under the
same policy, and \texttt{rlm\_query\_batched(prompts, context)} spawns a list.
Every \texttt{rlm\_query} invocation grows the trajectory tree $\mathcal{T}$ by
attaching a child whose own behaviour is again drawn from $\pi_\theta$.
Recursion turns the per-token causal-language-model MDP --- in which states are
prefixes and actions are next tokens \citep{schulman2017ppo} --- into a
tree-structured trajectory whose leaves all eventually call \texttt{FINAL}.

The motivation is long-context reasoning. Modern frontier models still struggle
when the relevant evidence is scattered across millions of tokens, e.g.\
multi-paper scientific corpora. The RLM scaffold lets a small model
\emph{decompose} a long-context query into many short-context child queries
(``search $\to$ expand $\to$ extract''), each of which fits comfortably in the
backbone's effective context. Chain-of-thought prompting \citep{wei2022cot}
unlocked sequential intermediate reasoning; the RLM idea is the natural next
agentic-scaling move, replacing a flat thought stream with a recursive tree.

\section*{Method}

The blog's central contribution is to RL-fine-tune a \emph{single} shared policy
$\pi_\theta$ to play both roles, rather than freezing the children as in the
original RLM setup. The optimisation primitive is GRPO \citep{shao2024deepseekmath,
deepseek2025r1}: for each query $x$, sample $G$ root rollouts
$\{y_g\}_{g=1}^G \sim \pi_{\theta_{\text{old}}}(\cdot \mid x)$, score each with a
rubric reward $r_g \in [0,1]$ (an LLM-judge call), and form group-relative
advantages
\[
A_g \;=\; \frac{r_g - \mu_r}{\sigma_r + \delta},
\qquad
\mu_r = \tfrac{1}{G}\textstyle\sum_g r_g,
\quad \sigma_r = \mathrm{stddev}(r_g).
\]

Each rollout (root or child) contributes a token-level PPO-clipped surrogate
\citep{schulman2017ppo}:
\begin{equation}
\mathcal{L}_{\text{node}}(y, A; \theta)
= \frac{1}{|y|} \sum_{t=1}^{|y|}
  \min\!\Big(\rho_\theta(y^{(t)}) A,\;
             \mathrm{clip}\!\big(\rho_\theta(y^{(t)}), 1{-}\varepsilon, 1{+}\varepsilon\big) A\Big).
\label{eq:node}
\end{equation}

The novel piece is the \textbf{advantage-inheritance} rule. Each child rollout
$y_{g,i}$ inherits the parent's GRPO advantage, $A_{g,i} := A_g$, and the
per-root contribution to the loss is balanced by $1/k_g$:
\begin{equation}
\mathcal{L}_g^{\text{full}}(\theta)
\;=\; \mathcal{L}_{\text{node}}(y_g, A_g; \theta)
   \;+\; \frac{1}{k_g} \sum_{i=1}^{k_g} \mathcal{L}_{\text{node}}(y_{g,i}, A_g; \theta).
\label{eq:depth1}
\end{equation}
Without the $1/k_g$ factor, a root that spawned more children would dominate
the gradient and the optimiser would develop a perverse incentive to call
\texttt{rlm\_query} more often. With it, every root contributes equally
regardless of how branchy its subtree is.

The depth-1 form \eqref{eq:depth1} extends recursively to arbitrary depth:
\begin{equation}
\mathcal{L}_{\text{subtree}}(y, A; \theta)
\;=\; \mathcal{L}_{\text{node}}(y, A; \theta)
   \;+\; \frac{1}{k_y} \sum_{i=1}^{k_y} \mathcal{L}_{\text{subtree}}(y_i, A; \theta).
\label{eq:subtree}
\end{equation}
Every descendant of root $g$ carries the same advantage $A_g$, and the
$1/k_y$ averaging at every level keeps the depth-balanced contribution
invariant under the local branching factor.

\section*{Results}

The Qwen3.5-4B backbone is first cold-started by SFT on teacher rollouts
distilled from Qwen3.5-397B-A17B, then RL-fine-tuned with the objective above.
Training runs on a single $8\times$H200 node with batch 16, $G=8$ root samples
per prompt and $k_g$ up to $4$, yielding up to $16 \cdot 8 \cdot 4 = 512$
concurrent rollouts per gradient step. The average rubric score on the training
distribution rises from $0.30$ to $0.60$ over RL.

On a held-out multi-paper evidence-selection benchmark, the RL-fine-tuned 4B
model reaches an average rubric score of $\approx 0.60$, matching Claude
Sonnet 4.6's $0.607$ on the same scaffold while running in $\approx 7$\,s per
query against Sonnet's $\approx 60$\,s --- an $8.5\times$ wall-clock speedup at
roughly two orders of magnitude fewer parameters. As a quality result this is a
\emph{tie}; the production win is latency and serving cost.

\section*{Position in the literature}

\paragraph{Original RLM paper.}
\citet{zhang2026rlm} introduced the recursive-LM scaffold and showed that a
\emph{prompted} (not RL-trained) frontier model can decompose long-context
queries into a tree of sub-rollouts. They trained only the root, leaving
children to a frozen LM. The blog reviewed here removes that asymmetry: one
policy, both roles, end-to-end RL.

\paragraph{Native parallelism.}
\citet{wu2025parallel} take a similar single-policy RL stance but replace the
recursion tree with a \emph{flat} parallel-reasoner ensemble distilled from the
model itself. The two papers share the ethos of training a small model to
exploit its own parallelism; they differ on whether the parallelism is
recursive (tree, this work) or flat (ensemble, theirs).

\paragraph{Pre-training / mid-training / RL interplay.}
\citet{zhang2025interplay} survey where in the modern training pipeline RL pays
off. Their findings frame why RLM training works at the 4B scale: a sufficiently
strong mid-trained backbone is a prerequisite for GRPO to lift rubric scores
the way the blog reports.

\paragraph{GRPO behaviour vs SFT.}
\citet{rajani2025grpo} characterise GRPO as a ``scalpel'' that amplifies
already-present capabilities while SFT acts as a ``hammer'' that overwrites
behaviour. The cold-start SFT $\to$ GRPO recipe used by the blog is exactly the
combination their analysis recommends: SFT to install the recursive scaffold
behaviour, GRPO to optimise it.

\paragraph{Rubric-based rewards.}
The blog uses an LLM-judge rubric reward instead of a verifiable reward (which
the authors report was too noisy: ``F1 of selected snippets'' had unusable
variance). This methodology comes from \citet{gunjal2025rubrics} and
\citet{yuan2025rubric}, both of whom argue that rubric rewards extend RL to
domains where ground-truth verifiers are absent --- exactly the
evidence-selection setting here.

\paragraph{Chain-of-thought.}
\citet{wei2022cot} unlocked sequential intermediate reasoning. Recursive LMs
are arguably the next agentic-scaling story: where CoT replaces a single token
prediction with a sequence of intermediate tokens, RLM replaces a single
context with a tree of contexts. The training-time challenge moves from
``elicit reasoning'' to ``credit-assign across a tree.''

\paragraph{DeepSeek-R1.}
\citet{deepseek2025r1} demonstrated GRPO-at-scale with rule-based rewards.
The blog inherits the GRPO machinery but swaps rule-based rewards for rubric
rewards and extends the credit-assignment to a recursive trajectory tree.

\paragraph{Hierarchical RL --- options framework.}
A connection the blog does not draw explicitly: RLM children are very nearly
\emph{options} in the sense of \citet{sutton1999options}. An option is a tuple
$(I, \pi_o, \beta)$ where $I$ is the initiation set, $\pi_o$ the option policy,
and $\beta$ the termination condition. For an RLM child, $I$ is the REPL state
at the moment of \texttt{rlm\_query}, $\pi_o = \pi_\theta$ (the shared
policy --- so options share parameters), and $\beta$ fires precisely when the
child calls \texttt{FINAL}. Read this way, the recursive-subtree loss
\eqref{eq:subtree} is a temporally-abstract policy gradient that backs up the
\emph{root-level} advantage through every option in the tree; the $1/k_g$
averaging is what re-weights options so their density does not bias the root
gradient. Importing convergence and unbiasedness theorems from the options
literature is, in our view, the cleanest theoretical hardening of the blog's
informal unbiasedness argument.

\paragraph{Engineering substrate.}
The implementation lives in \citet{skyrl2025} (specifically PR \#1596), an
extension of the SkyRL framework for recursive and agentic LM training.

\section*{Discussion}

\paragraph{Well-supported.}
The core empirical claim --- that a single shared policy with parent-advantage
inheritance and $1/k_g$ averaging trains a 4B RLM to Sonnet-4.6 parity on
evidence selection --- is well-supported by the reported numbers. The
$1/k_g$ balancing and the cold-start SFT recipe are both standard moves in the
GRPO literature \citep{shao2024deepseekmath, deepseek2025r1, rajani2025grpo}
and inherit those papers' training-stability story.

\paragraph{Hand-waved.}
The unbiasedness of advantage inheritance is argued informally. Inheriting
$A_g$ for every descendant is unbiased only when the child's reward
contribution is $\theta$-independent given the parent's rollout window; this
holds for evidence selection (parent reward is roughly additive in child
contributions) but fails for tasks with non-additive child interactions
(e.g.\ tool-use chains where child A's output is child B's input). The blog
flags ``more fine-grained credit assignment'' as future work but does not
formalise the load-bearing assumption. Likewise, the model still slightly
trails Sonnet 4.6 on rubric quality ($0.60$ vs $0.607$); the headline result
is wall-clock parity, not quality breakthrough.

\paragraph{Future directions.}
Three open questions, expanded in the companion \texttt{05-improvements.tex}:
(i) does the recursive-subtree loss \eqref{eq:subtree} train at depth $d > 1$?
The blog only trains depth-1; (ii) at fixed compute, what is the optimal $k_g$,
and is there a knee in the rubric-vs-$k_g$ curve? (iii) can RL discover the
search-expand-extract strategy from scratch rather than being prompted into it?

\section*{See also: cross-links to the \texttt{first-principles-to-llms} chain}

This paper sits naturally on top of several chapters of the chain repository
\href{https://github.com/pleyva2004/first-principles-to-llms}{\texttt{first-principles-to-llms}}:

\begin{itemize}
\item Causal-LM-as-MDP framing
  (\href{https://github.com/pleyva2004/first-principles-to-llms/tree/main/Chapter\%2025}{Ch~25})
  --- recovered at every node of the trajectory tree.
\item Tiny-GPT pre-training
  (\href{https://github.com/pleyva2004/first-principles-to-llms/tree/main/Chapter\%2027}{Ch~27})
  --- the backbone scale at which a GRPO-trained model can match a much
  larger frontier model.
\item SFT, RLHF, DPO
  (\href{https://github.com/pleyva2004/first-principles-to-llms/tree/main/Chapter\%2028}{Ch~28})
  --- the cold-start SFT step that precedes RL here.
\item MDP foundations
  (\href{https://github.com/pleyva2004/first-principles-to-llms/tree/main/Chapter\%2029}{Ch~29})
  --- where the per-token MDP and value/advantage definitions live.
\item Maximum-entropy RL
  (\href{https://github.com/pleyva2004/first-principles-to-llms/tree/main/Chapter\%2030}{Ch~30})
  --- conceptual prior for the PPO clipping/KL machinery.
\item Policy gradient and the GRPO derivation
  (\href{https://github.com/pleyva2004/first-principles-to-llms/tree/main/Chapter\%2031}{Ch~31})
  --- the surrogate \eqref{eq:node} and the group-baseline variance-reduction
  argument that the blog inherits.
\end{itemize}

\bibliographystyle{plainnat}
\bibliography{references}

\end{document}
